%% 第三章 无穷级数的解析解研究（约4000字）

\chapter{无穷级数的解析解研究}

本章在第二章理论基础之上，系统研究无穷级数解析解的存在条件与主要求解方法，并对解析解的局限性进行分析。解析解以有限的数学符号精确到达无穷过程之和，是级数求解的理想结果。然而，并非所有收敛级数都拥有解析解，认识这一边界是构建谱系化框架的前提。

\section{解析解的存在条件}

解析解的存在性并非普遍成立，它取决于级数的结构类型与相应的理论保证。不同类型的无穷级数有着各自的解析解存在条件，本节分别就常数项级数、幂级数和傅里叶级数展开讨论。

\subsection{常数项级数的解析解条件}

常数项级数 $\displaystyle \sum_{n=1}^{\infty} a_n$ 具有解析解，意味着其和 $S$ 可以由初等函数的闭合表达式精确表示。这一条件并无统一的充要判定准则，通常需要结合级数的具体结构来判断。

能够获得解析解的常数项级数，往往具有以下典型结构特征：

\textbf{（1）等比结构。}若级数各项之比为常数，即 $a_{n+1}/a_n = r$（$\abs{r} < 1$），则级数为等比级数，其和 $S = a/(1-r)$（$a$ 为首项）为初等闭合表达式。

\textbf{（2）裂项结构。}若通项 $a_n$ 可以分解为形如 $f(n) - f(n+1)$ 的差形式，则部分和可发生大量抵消，最终只剩有限项，和为初等表达式。

\textbf{（3）已知函数在特定点的展开。}若级数恰好是某已知初等函数在某一点的幂级数展开（如 $e^x$、$\sin x$、$\ln(1+x)$ 等），则其和等于该函数在该点的值，即为解析解。

\textbf{（4）特殊结构常数。}部分常数项级数的和虽非有理数，但可以用著名数学常数（如 $\pi$、$e$、$\ln 2$ 等）的初等组合表达，同样视为解析解。典型例子为巴塞尔问题
\begin{equation}
\sum_{n=1}^{\infty} \frac{1}{n^2} = \frac{\pi^2}{6}
\end{equation}
其证明需要借助欧拉对正弦函数无穷乘积展开的深刻分析。

反之，若级数的通项既不具备上述结构，又无法与已知函数建立联系，则其和通常不存在初等闭合表达式，例如 $\displaystyle \sum_{n=1}^{\infty} 1/n^3$（阿佩里常数 $\zeta(3)$）至今没有已知的解析表达式。

\subsection{幂级数的解析解条件}

幂级数 $\displaystyle \sum_{n=0}^{\infty} c_n (x - x_0)^n$ 在其收敛半径内定义了一个函数，而该函数是否具有初等解析表达式，取决于其系数序列 $\{c_n\}$ 是否与某已知初等函数的泰勒展开系数相吻合。

\begin{theorem}[幂级数解析性条件]
若幂级数 $\displaystyle \sum_{n=0}^{\infty} c_n (x - x_0)^n$ 的系数满足
\begin{equation}
c_n = \frac{f^{(n)}(x_0)}{n!}
\end{equation}
其中 $f(x)$ 是某个在 $x_0$ 处解析的初等函数，且 $f$ 的泰勒级数在某一邻域内处处收敛于 $f$ 本身，则该幂级数在收敛域内的和函数等于 $f(x)$，从而具有解析表达式。
\end{theorem}

这一条件实质上要求幂级数"恰好来自"某个初等函数的泰勒展开。常见的具有解析解的幂级数包括：
\begin{enumerate}[label=（\arabic*）]
  \item 指数函数展开：$\displaystyle \sum_{n=0}^{\infty} x^n/n! = e^x$，收敛域为 $(-\infty, +\infty)$；
  \item 正弦函数展开：$\displaystyle \sum_{n=0}^{\infty} (-1)^n x^{2n+1}/(2n+1)! = \sin x$，收敛域为 $(-\infty, +\infty)$；
  \item 对数函数展开：$\displaystyle \sum_{n=1}^{\infty} (-1)^{n+1} x^n/n = \ln(1+x)$，收敛域为 $(-1, 1]$；
  \item 等比级数展开：$\displaystyle \sum_{n=0}^{\infty} x^n = 1/(1-x)$，收敛域为 $(-1, 1)$。
\end{enumerate}

需要特别说明的是，并非所有幂级数都与初等函数对应。例如，误差函数
\begin{equation}
\operatorname{erf}(x) = \frac{2}{\sqrt{\pi}}\int_0^x e^{-t^2}\, \mathrm{d}t = \frac{2}{\sqrt{\pi}} \sum_{n=0}^{\infty} \frac{(-1)^n x^{2n+1}}{n!(2n+1)}
\end{equation}
虽然收敛域为全实数轴，但其和函数 $\operatorname{erf}(x)$ 本身不是初等函数，因此该幂级数的和不具有初等解析表达式，只能以特殊函数的形式记录。

\subsection{傅里叶级数的解析解条件}

傅里叶级数 $\displaystyle a_0/2 + \sum_{n=1}^{\infty} (a_n \cos nx + b_n \sin nx)$ 的解析性问题与其收敛性以及和函数的表达密切相关。傅里叶级数的收敛定理给出了存在解析解的基本条件：

\begin{theorem}[狄利克雷收敛定理]
设 $f(x)$ 是以 $2\pi$ 为周期的函数，若 $f(x)$ 在一个周期内满足狄利克雷条件（即在一个周期上绝对可积、分段单调且只有有限个第一类间断点），则 $f(x)$ 的傅里叶级数在每个连续点 $x$ 处收敛于 $f(x)$，在每个间断点 $x_0$ 处收敛于
\begin{equation}
\frac{f(x_0^+) + f(x_0^-)}{2}.
\end{equation}
\end{theorem}

满足狄利克雷条件的函数的傅里叶级数，其和函数可用 $f(x)$ 本身（或其左右极限的均值）来表达，这就构成了傅里叶级数的解析解。当 $f(x)$ 本身是初等函数时，傅里叶级数的和也可写成初等表达式。

然而，傅里叶级数即使在某点处收敛，在端点或间断点处的行为也可能出现吉布斯现象——即部分和在间断点附近产生约 $9\%$ 的超调现象\upcite{stein2003fourier}。这并不影响级数在连续点的收敛性，但提示了傅里叶展开的局部误差特性。

\section{主要求解方法与实例}

当确认无穷级数具有解析解后，关键问题是如何系统地求出这一解。本节介绍四种主要的解析求解方法，每种方法均配以典型实例加以说明。

\subsection{方法一：等比级数求和公式}

等比级数是最基本也是应用最广泛的具有解析解的级数类型。其求和依据为等比数列的部分和公式的极限形式。

\begin{theorem}[等比级数求和定理]
设公比为 $r$，首项为 $a$（$a \neq 0$），则等比级数
\begin{equation}
\sum_{n=0}^{\infty} a r^n = a + ar + ar^2 + \cdots
\end{equation}
当 $\abs{r} < 1$ 时收敛，其和为
\begin{equation} \label{eq:geom}
S = \frac{a}{1-r};
\end{equation}
当 $\abs{r} \geq 1$ 时级数发散。
\end{theorem}

\begin{proof}
第 $n$ 个部分和为
\begin{equation}
S_n = a \cdot \frac{1 - r^n}{1 - r} \quad (r \neq 1).
\end{equation}
当 $\abs{r} < 1$ 时，$r^n \to 0$，故 $S_n \to a/(1-r)$。
\end{proof}

等比级数公式的重要性不仅在于其自身的求和，还在于它是许多其他级数求解的基础。通过变量替换、逐项求导或逐项积分，可以从等比级数出发推导出多种级数的解析解。

\begin{example}[含参数的等比级数变形]
求级数 $\displaystyle \sum_{n=1}^{\infty} n x^{n-1}$（$\abs{x} < 1$）的和。

\textbf{解：}注意到 $\displaystyle \sum_{n=0}^{\infty} x^n = 1/(1-x)$（$\abs{x} < 1$）。由于幂级数在收敛域内可以逐项求导，对两边关于 $x$ 求导得
\begin{equation}
\sum_{n=1}^{\infty} n x^{n-1} = \frac{\mathrm{d}}{\mathrm{d}x} \cdot \frac{1}{1-x} = \frac{1}{(1-x)^2}.
\end{equation}
故所求级数的和函数为 $S(x) = 1/(1-x)^2$，$\abs{x} < 1$。
\end{example}

\begin{example}[等比级数的积分变形]
求级数 $\displaystyle \sum_{n=0}^{\infty} (-1)^n x^{2n+1}/(2n+1)$（$\abs{x} \leq 1$）的和。

\textbf{解：}设 $\displaystyle S(x) = \sum_{n=0}^{\infty} (-1)^n x^{2n+1}/(2n+1)$，则
\begin{equation}
S'(x) = \sum_{n=0}^{\infty} (-1)^n x^{2n} = \frac{1}{1+x^2} \quad (\abs{x} < 1).
\end{equation}
积分得 $S(x) = \arctan x + C$。由 $S(0) = 0$ 知 $C = 0$，故 $S(x) = \arctan x$（$\abs{x} < 1$）。对于端点 $x = 1$，原级数满足莱布尼兹条件，条件收敛；由阿贝尔定理，$S(x)$ 在 $x \to 1^-$ 时的极限等于级数在 $x=1$ 处的和，故
\begin{equation}
\sum_{n=0}^{\infty} \frac{(-1)^n}{2n+1} = \arctan 1 = \frac{\pi}{4}.
\end{equation}
这即为著名的莱布尼兹公式。
\end{example}

\subsection{方法二：裂项求和}

裂项求和法的核心思想是将通项 $a_n$ 表示为 $f(n) - f(n+1)$ 的形式，使得部分和大量抵消，只保留首尾两项，从而直接得到闭合表达式：
\begin{equation}
S_n = \sum_{k=1}^{n} [f(k) - f(k+1)] = f(1) - f(n+1).
\end{equation}
若 $\displaystyle \lim_{n \to \infty} f(n+1) = L$ 存在，则级数收敛，和为 $S = f(1) - L$。

\begin{example}[有理分式裂项]
求级数 $\displaystyle\sum_{n=1}^{\infty} \frac{1}{n(n+1)}$ 的和。

\textbf{解：}利用部分分数分解，
\begin{equation}
\frac{1}{n(n+1)} = \frac{1}{n} - \frac{1}{n+1}.
\end{equation}
部分和为
\begin{equation}
S_n = \sum_{k=1}^{n}\left(\frac{1}{k} - \frac{1}{k+1}\right) = 1 - \frac{1}{n+1}.
\end{equation}
令 $n \to \infty$，得 $S = 1$。
\end{example}

\begin{example}[含平方的裂项]
求级数 $\displaystyle\sum_{n=1}^{\infty} \frac{1}{n(n+2)}$ 的和。

\textbf{解：}分解为
\begin{equation}
\frac{1}{n(n+2)} = \frac{1}{2}\left(\frac{1}{n} - \frac{1}{n+2}\right).
\end{equation}
部分和为
\begin{equation}
S_n = \frac{1}{2}\left(1 + \frac{1}{2} - \frac{1}{n+1} - \frac{1}{n+2}\right).
\end{equation}
令 $n \to \infty$，得 $S = (1/2) \cdot (3/2) = 3/4$。
\end{example}

裂项求和法适用于通项可进行部分分数分解、或具有差分结构的级数，在有理函数项级数中尤为有效。其局限在于对通项结构要求较高，不具备普适性。

\subsection{方法三：泰勒展开法}

泰勒展开法是将已知函数展开为幂级数，从而将特定数值级数与函数在某点的取值建立联系，进而获得解析解。

设 $f(x)$ 在 $x_0$ 处无穷次可微，则其泰勒展开为
\begin{equation}
f(x) = \sum_{n=0}^{\infty} \frac{f^{(n)}(x_0)}{n!} (x - x_0)^n.
\end{equation}
当 $x_0 = 0$ 时，称为麦克劳林展开。常用的麦克劳林展开包括（收敛域各有不同）：
\begin{align}
e^x &= \sum_{n=0}^{\infty} \frac{x^n}{n!}, \quad x \in (-\infty, +\infty) \label{eq:exp}\\
\sin x &= \sum_{n=0}^{\infty} \frac{(-1)^n x^{2n+1}}{(2n+1)!}, \quad x \in (-\infty, +\infty) \notag\\
\cos x &= \sum_{n=0}^{\infty} \frac{(-1)^n x^{2n}}{(2n)!}, \quad x \in (-\infty, +\infty) \notag\\
\ln(1+x) &= \sum_{n=1}^{\infty} \frac{(-1)^{n+1} x^n}{n}, \quad x \in (-1, 1] \label{eq:ln}\\
(1+x)^\alpha &= \sum_{n=0}^{\infty} \binom{\alpha}{n} x^n, \quad x \in (-1, 1) \notag
\end{align}

\begin{example}[利用指数函数展开]
求级数 $\displaystyle\sum_{n=0}^{\infty} \frac{(-1)^n}{n!}$ 的和。

\textbf{解：}由公式\eqref{eq:exp}，令 $x = -1$，得
\begin{equation}
\sum_{n=0}^{\infty} \frac{(-1)^n}{n!} = e^{-1} = \frac{1}{e}.
\end{equation}
\end{example}

\begin{example}[利用对数展开与特殊值]
求级数 $\displaystyle\sum_{n=1}^{\infty} \frac{(-1)^{n+1}}{n}$ 的和。

\textbf{解：}由公式\eqref{eq:ln}，令 $x = 1$（此时级数在 $x = 1$ 处条件收敛，由阿贝尔定理可在端点取值），得
\begin{equation}
\sum_{n=1}^{\infty} \frac{(-1)^{n+1}}{n} = \ln 2.
\end{equation}
这一结果表明交错调和级数具有精确的解析解。
\end{example}

\begin{example}[逐项积分获得解析解]
求级数 $\displaystyle\sum_{n=0}^{\infty} \frac{x^{n+1}}{n+1}$（$\abs{x} < 1$）的和。

\textbf{解：}设和函数为 $S(x)$，则 $\displaystyle S'(x) = \sum_{n=0}^{\infty} x^n = 1/(1-x)$。由 $S(0) = 0$，积分得
\begin{equation}
S(x) = -\ln(1-x), \quad \abs{x} < 1.
\end{equation}
\end{example}

泰勒展开法的优势在于充分利用了初等函数的已知展开结果，适用范围较广，对于含阶乘、幂函数的系数级数往往能与初等函数建立联系。

\subsection{方法四：傅里叶级数三角展开}

傅里叶展开法是将特定周期函数展开为三角级数，并利用狄利克雷定理在特定点取值，从而将一个看似无关的数值级数与函数值联系起来，获得精确解析解。

\begin{example}[利用傅里叶展开求 $\zeta(2)$]
求巴塞尔级数 $\displaystyle\sum_{n=1}^{\infty} \frac{1}{n^2}$ 的和。

\textbf{解：}考虑 $f(x) = x^2$ 在 $[-\pi, \pi]$ 上的傅里叶展开。计算傅里叶系数：
\begin{align}
a_0 &= \frac{1}{\pi}\int_{-\pi}^{\pi} x^2 \, \mathrm{d}x = \frac{2\pi^2}{3}, \notag\\
a_n &= \frac{1}{\pi}\int_{-\pi}^{\pi} x^2 \cos nx \, \mathrm{d}x = \frac{4(-1)^n}{n^2}, \quad n \geq 1, \notag\\
b_n &= 0 \quad \text{（$f$ 为偶函数）.} \notag
\end{align}
故傅里叶展开为
\begin{equation}
x^2 = \frac{\pi^2}{3} + \sum_{n=1}^{\infty} \frac{4(-1)^n}{n^2} \cos nx.
\end{equation}
令 $x = \pi$（连续点，级数收敛于 $f(\pi) = \pi^2$），得
\begin{equation}
\pi^2 = \frac{\pi^2}{3} + 4\sum_{n=1}^{\infty} \frac{1}{n^2}.
\end{equation}
整理得
\begin{equation}
\sum_{n=1}^{\infty} \frac{1}{n^2} = \frac{\pi^2}{6}.
\end{equation}
\end{example}

这一方法将正整数倒数平方的无穷和与超越常数 $\pi^2$ 相联系，体现了解析解背后的数学结构。

\begin{example}[锯齿波展开]
设 $f(x) = x$（$-\pi < x < \pi$），以 $2\pi$ 为周期延拓。其傅里叶展开为
\begin{equation}
x = 2\sum_{n=1}^{\infty} \frac{(-1)^{n+1}}{n} \sin nx, \quad x \in (-\pi, \pi).
\end{equation}
令 $x = \pi/2$，由 $\sin(n\pi/2)$ 的规律可得
\begin{equation}
\frac{\pi}{4} = 1 - \frac{1}{3} + \frac{1}{5} - \frac{1}{7} + \cdots = \sum_{n=0}^{\infty} \frac{(-1)^n}{2n+1}.
\end{equation}
此即莱布尼兹-格雷戈里公式。
\end{example}

傅里叶展开法的特点在于，它将计算傅里叶系数的积分操作与利用狄利克雷定理的级数求和操作结合起来，使得许多用代数方法难以处理的数值级数获得了精确的解析表达。

\section{解析解的局限性分析}

尽管解析解具有精确性和理论价值，但其适用范围受到严格限制。理解这些局限性，对于后续引入数值解方法以及构建谱系化框架是必要的。

\subsection{存在性局限}

解析解的最根本局限在于其存在性不具有普遍性。大量收敛的无穷级数，其和不能用初等函数的闭合表达式来表示。

一方面，从分析学的超越性理论可知，大多数分析函数（如 Gamma 函数、Bessel 函数、误差函数等）并非初等函数，因此与这些函数相关的级数求和自然不具有初等解析解。例如：
\begin{enumerate}[label=（\arabic*）]
  \item $\displaystyle\sum_{n=0}^{\infty} \frac{1}{(n!)^2}$ 的和为 $I_0(2)$，其中 $I_0$ 是第一类零阶修正贝塞尔函数，不能用初等函数表达；
  \item $\displaystyle\sum_{n=1}^{\infty} \frac{1}{n^3}$（即黎曼 $\zeta$ 函数值 $\zeta(3)$，阿佩里常数）虽已被证明为无理数\upcite{apery1979irrationalite}，但至今没有已知的初等闭合表达式；
  \item $\displaystyle\sum_{n=1}^{\infty} \frac{1}{n^5}$、$\displaystyle\sum_{n=1}^{\infty} \frac{1}{n^7}$ 等奇数次幂 $\zeta$ 函数值的闭合形式同样未知。
\end{enumerate}

另一方面，即使级数收敛，其收敛速度极慢，在有限计算步骤内实际上不可能通过部分和得到高精度近似，这从实用角度进一步限制了解析解的可获得性。

\subsection{运算局限}

即便某个级数在形式上具有解析解，求解过程本身也可能面临运算上的障碍。

\textbf{（1）积分困难。}对于由已知函数的幂级数展开推导解析解的方法，常需要计算相关的不定积分。若被积函数本身没有初等原函数（如 $\displaystyle \int e^{-x^2}\, \mathrm{d}x$、$\displaystyle \int (\sin x)/x\, \mathrm{d}x$ 等），则通过积分手段求解析解的路径便行不通。

\textbf{（2）傅里叶系数的可计算性。}傅里叶展开法要求计算傅里叶系数，这等价于计算特定的定积分。若函数形式复杂，积分无法用初等函数表达，则该方法同样受阻。

\textbf{（3）部分分数分解的复杂度。}裂项求和法在面对高阶或多因子的有理函数分母时，部分分数分解的计算量迅速增大，且并非总能化简为简洁形式。对于无理函数或超越函数通项的级数，裂项结构往往不存在。

\textbf{（4）条件收敛的精细性。}对于条件收敛级数（如交错调和级数），虽然可以求得解析解，但由黎曼重排定理知其项的重排可改变和值。这要求在利用级数的代数运算（如逐项微积分）时必须验证相关条件，稍有疏忽便可能导致错误结论。

\subsection{收敛域局限}

解析解通常只在有限的区域内有效，这构成了解析方法的第三类局限。

对于幂级数 $\displaystyle \sum_{n=0}^{\infty} c_n x^n$，即使其和函数 $f(x)$ 在整个实数轴上有意义（如 $e^x$、$\sin x$），幂级数展开的使用范围仍限于其收敛域 $\abs{x} < R$（或 $\abs{x} \leq R$ 的部分端点）。对于收敛半径有限的幂级数（如 $\displaystyle \sum_{n=0}^{\infty} x^n = 1/(1-x)$，$R = 1$），其解析解只在 $\abs{x} < 1$ 时成立，在 $\abs{x} \geq 1$ 时级数本身已发散，解析表达式无法延伸。

对于傅里叶级数，解析解在连续点处成立，但在间断点处级数收敛于左右极限的均值，而非函数值本身。在实际应用中，若函数有间断点，解析解的适用范围需要额外的说明和处理。

此外，在多变量推广的情形中，幂级数的收敛域从实数轴上的区间变为复平面中的圆盘（或更复杂的区域），单变量解析解的形式不能直接套用。

综合上述三类局限可以看出：解析解在\textbf{精确性}上具有优势，但其存在性不具普遍保证，运算过程受初等函数体系的制约，有效范围也局限于收敛域之内。这些局限使得数值方法在无穷级数求解中不可或缺。若能将两者纳入统一框架，明确各自的适用边界与互补关系，将有助于提高对无穷级数求解策略的系统性认识。这一框架的具体构建，将在第五章展开。
